% GENERATED FILE. Edit canonical lessons, then run npm run generate:books.
% canonical-source-hash: fnv1a64:318da2bae937635c
% canonical-lessons: GU-C471-van, GU-C471-kad, GU-C471-deh, GU-C471-valiya, GU-C471-paliya

\chapter{How People Look}
\label{ch:gu-how-people-look}
\hlchaptermodality{471}

\begin{chapteropening}
\textbf{By the end of this chapter:} \emph{I can describe how someone looks.}

Complete the last lesson of chapter 471: I can describe how someone looks.
\end{chapteropening}

\section[\texorpdfstring{\gu{વાન} (vān)}{વાન (vān)}]{\gu{વાન} (vān) — complexion; skin colour}

\label{lesson:GU-C471-van}



\begin{quote}
Before the new one: say the Gujarati for a lock of hair, then the Gujarati for the parting of the hair.
\end{quote}

\subsection*{You'll want to know: \gu{વાન}}
\textbf{\gu{વાન}} — \emph{vān} — \textquotedblleft{}complexion; skin colour\textquotedblright{}.

\begin{quote}
\gu{દાદા} \gu{કહે} \gu{છે} \gu{કે} \gu{મારો} \gu{વાન} \gu{મારી} \gu{માતાનો} \gu{છે}.
\end{quote}

\begin{quote}
\emph{dādā kahe che ke māro vān mārī mātāno che.}
\end{quote}

\begin{quote}
Grandfather says that I have my mother's complexion.
\end{quote}

A masculine noun.

\begin{cousinweb}
\textbf{\gu{વાન}} (\emph{vān}) is inherited from Sanskrit \emph{varṇa}, \textquotedblleft{}appearance, colour\textquotedblright{}; CDIAL lists the Gujarati word, \textquotedblleft{}colour\textquotedblright{} (CDIAL 11338).
\end{cousinweb}

\begin{grammarlens}[title={What you've built}]
One more word for describing what you lived and felt.
\end{grammarlens}

\subsection*{Your turn}
\begin{itemize}
\raggedright
  \item \emph{Say it:} \emph{vān}
  \item \emph{Say it:} \emph{vān}, once more
  \item \emph{Say it:} \emph{dādā kahe che ke māro vān mārī mātāno che.} — Grandfather says that I have my mother's complexion.
  \item \emph{Recall:} say the Gujarati for a lock of hair, then the Gujarati for the parting of the hair, then say \emph{vān} again
\end{itemize}

\begin{tcolorbox}[breakable,colback=teal!4,colframe=teal!35!black,title={Before you move on}]
What does \gu{વાન} mean? (\textquotedblleft{}Complexion; skin colour\textquotedblright{}.) Say it once more.
\end{tcolorbox}
\section[\texorpdfstring{\gu{કદ} (kad)}{કદ (kad)}]{\gu{કદ} (kad) — height; size (of a person)}

\label{lesson:GU-C471-kad}



\begin{quote}
Before the new one: say the Gujarati for the parting of the hair, then the Gujarati for complexion.
\end{quote}

\subsection*{You'll want to know: \gu{કદ}}
\textbf{\gu{કદ}} — \emph{kad} — \textquotedblleft{}height; size (of a person)\textquotedblright{}.

\begin{quote}
\gu{અરુણનું} \gu{કદ} \gu{પિતા} \gu{કરતાં} \gu{મોટું} \gu{છે}.
\end{quote}

\begin{quote}
\emph{Aruṇnũ kad pitā kartā̃ moṭũ che.}
\end{quote}

\begin{quote}
Arun is taller than his father.
\end{quote}

A neuter noun.

\begin{cousinweb}
\textbf{\gu{કદ}} (\emph{kad}) is usually traced to Arabic \emph{qadd}, \textquotedblleft{}stature\textquotedblright{}, which reached Gujarati through Persian.
\end{cousinweb}

\begin{grammarlens}[title={What you've built}]
One more word for describing what you lived and felt.
\end{grammarlens}

\subsection*{Your turn}
\begin{itemize}
\raggedright
  \item \emph{Say it:} \emph{kad}
  \item \emph{Say it:} \emph{kad}, once more
  \item \emph{Say it:} \emph{Aruṇnũ kad pitā kartā̃ moṭũ che.} — Arun is taller than his father.
  \item \emph{Recall:} say the Gujarati for the parting of the hair, then the Gujarati for complexion, then say \emph{kad} again
\end{itemize}

\begin{tcolorbox}[breakable,colback=teal!4,colframe=teal!35!black,title={Before you move on}]
What does \gu{કદ} mean? (\textquotedblleft{}Height; size (of a person)\textquotedblright{}.) Say it once more.
\end{tcolorbox}
\section[\texorpdfstring{\gu{દેહ} (deh)}{દેહ (deh)}]{\gu{દેહ} (deh) — the body (in careful speech); build}

\label{lesson:GU-C471-deh}



\begin{quote}
Before the new one: say the Gujarati for complexion, then the Gujarati for height.
\end{quote}

\subsection*{You'll want to know: \gu{દેહ}}
\textbf{\gu{દેહ}} — \emph{deh} — \textquotedblleft{}the body (in careful speech); build\textquotedblright{}.

\begin{quote}
\gu{કસરતને} \gu{લીધે} \gu{અરુણનો} \gu{દેહ} \gu{મજબૂત} \gu{છે}.
\end{quote}

\begin{quote}
\emph{kasratne līdhe Aruṇno deh majbūt che.}
\end{quote}

\begin{quote}
Thanks to exercise, Arun has a strong build.
\end{quote}

A masculine noun.

\begin{cousinweb}
\textbf{\gu{દેહ}} (\emph{deh}) is a tatsama, Sanskrit \emph{deha}, \textquotedblleft{}the body\textquotedblright{}, as Monier-Williams records it, from a root meaning \textquotedblleft{}to mould, to fashion\textquotedblright{}; everyday Gujarati says \textbf{\gu{શરીર}} (\emph{śarīr}).
\end{cousinweb}

\begin{grammarlens}[title={What you've built}]
One more word for describing what you lived and felt.
\end{grammarlens}

\subsection*{Your turn}
\begin{itemize}
\raggedright
  \item \emph{Say it:} \emph{deh}
  \item \emph{Say it:} \emph{deh}, once more
  \item \emph{Say it:} \emph{kasratne līdhe Aruṇno deh majbūt che.} — Thanks to exercise, Arun has a strong build.
  \item \emph{Recall:} say the Gujarati for complexion, then the Gujarati for height, then say \emph{deh} again
\end{itemize}

\begin{tcolorbox}[breakable,colback=teal!4,colframe=teal!35!black,title={Before you move on}]
What does \gu{દેહ} mean? (\textquotedblleft{}The body (in careful speech); build\textquotedblright{}.) Say it once more.
\end{tcolorbox}
\section[\texorpdfstring{\gu{વળિયાં} (vaḷiyā̃)}{વળિયાં (vaḷiyā̃)}]{\gu{વળિયાં} (vaḷiyā̃) — wrinkles}

\label{lesson:GU-C471-valiya}



\begin{quote}
Before the new one: say the Gujarati for height, then the Gujarati for the body.
\end{quote}

\subsection*{You'll want to know: \gu{વળિયાં}}
\textbf{\gu{વળિયાં}} — \emph{vaḷiyā̃} — \textquotedblleft{}wrinkles\textquotedblright{}.

\begin{quote}
\gu{દાદીના} \gu{ચહેરા} \gu{પર} \gu{વળિયાં} \gu{છે}.
\end{quote}

\begin{quote}
\emph{dādīnā caherā par vaḷiyā̃ che.}
\end{quote}

\begin{quote}
There are wrinkles on my grandmother's face.
\end{quote}

A neuter noun, used in the plural.

\begin{cousinweb}
\textbf{\gu{વળિયાં}} (\emph{vaḷiyā̃}) is inherited from Sanskrit \emph{vali}, \textquotedblleft{}a fold of skin\textquotedblright{}; CDIAL lists the Gujarati word as \textquotedblleft{}wrinkles in a garment\textquotedblright{} (CDIAL 11412), and Gujarati uses it for the folds of the skin too.
\end{cousinweb}

\begin{grammarlens}[title={What you've built}]
One more word for describing what you lived and felt.
\end{grammarlens}

\subsection*{Your turn}
\begin{itemize}
\raggedright
  \item \emph{Say it:} \emph{vaḷiyā̃}
  \item \emph{Say it:} \emph{vaḷiyā̃}, once more
  \item \emph{Say it:} \emph{dādīnā caherā par vaḷiyā̃ che.} — There are wrinkles on my grandmother's face.
  \item \emph{Recall:} say the Gujarati for height, then the Gujarati for the body, then say \emph{vaḷiyā̃} again
\end{itemize}

\begin{tcolorbox}[breakable,colback=teal!4,colframe=teal!35!black,title={Before you move on}]
What does \gu{વળિયાં} mean? (\textquotedblleft{}Wrinkles\textquotedblright{}.) Say it once more.
\end{tcolorbox}
\section[\texorpdfstring{\gu{પળિયાં} (paḷiyā̃)}{પળિયાં (paḷiyā̃)}]{\gu{પળિયાં} (paḷiyā̃) — grey hairs}

\label{lesson:GU-C471-paliya}



\begin{quote}
Before the new one: say the Gujarati for the body, then the Gujarati for wrinkles.
\end{quote}

\subsection*{You'll want to know: \gu{પળિયાં}}
\textbf{\gu{પળિયાં}} — \emph{paḷiyā̃} — \textquotedblleft{}grey hairs\textquotedblright{}.

\begin{quote}
\gu{પિતાના} \gu{વાળમાં} \gu{હવે} \gu{પળિયાં} \gu{આવ્યાં} \gu{છે}.
\end{quote}

\begin{quote}
\emph{pitānā vāḷmā̃ have paḷiyā̃ āvyā̃ che.}
\end{quote}

\begin{quote}
Father's hair has started to turn grey.
\end{quote}

A neuter noun, used in the plural.

\begin{cousinweb}
\textbf{\gu{પળિયાં}} (\emph{paḷiyā̃}) is inherited from Sanskrit \emph{palita}, \textquotedblleft{}grey, hoary, old\textquotedblright{}; CDIAL lists the Gujarati word, \textquotedblleft{}grey hair, old age\textquotedblright{} (CDIAL 7962).
\end{cousinweb}

\begin{grammarlens}[title={What you've built}]
One more word for describing what you lived and felt.
\end{grammarlens}

\subsection*{Your turn}
\begin{itemize}
\raggedright
  \item \emph{Say it:} \emph{paḷiyā̃}
  \item \emph{Say it:} \emph{paḷiyā̃}, once more
  \item \emph{Say it:} \emph{pitānā vāḷmā̃ have paḷiyā̃ āvyā̃ che.} — Father's hair has started to turn grey.
  \item \emph{Recall:} say the Gujarati for the body, then the Gujarati for wrinkles, then say \emph{paḷiyā̃} again
\end{itemize}

\begin{tcolorbox}[breakable,colback=teal!4,colframe=teal!35!black,title={Before you move on}]
What does \gu{પળિયાં} mean? (\textquotedblleft{}Grey hairs\textquotedblright{}.) Say it once more.
\end{tcolorbox}
